% Final LaTeX version generated from the reviewed manuscript.
% Compile with the bundled Tectonic runtime or XeLaTeX + BibTeX.
\documentclass[11pt,a4paper]{article}

\usepackage[margin=0.79in]{geometry}
\usepackage{lmodern}
\usepackage[english]{babel}
\usepackage{microtype}
\usepackage{amsmath,amssymb}
\usepackage{booktabs,longtable,array,tabularx,makecell}
\usepackage{calc}
\usepackage{enumitem}
\usepackage{caption}
\usepackage{float}
\usepackage[numbers,sort&compress]{natbib}
\usepackage{hyperref}
\usepackage{xurl}
\usepackage{ragged2e}
\usepackage{titlesec}
\newcounter{none} % keeps Pandoc's unnumbered longtable wrapper valid

\hypersetup{
  colorlinks=true,
  linkcolor=black,
  citecolor=black,
  urlcolor=blue,
  breaklinks=true,
  pdftitle={Development of a Multitask Transformer-Based Model and Web System for Joint Sentiment Analysis and AI-Generated Review Provenance Classification},
  pdfauthor={Dias Kazikhanov and Aigul Meirmanova}
}

\setlength{\parindent}{0.5in}
\setlength{\parskip}{0pt}
\setlength{\emergencystretch}{3em}
\setlength{\tabcolsep}{4pt}
\newcolumntype{L}[1]{>{\raggedright\arraybackslash}p{#1}}
\newcolumntype{C}[1]{>{\centering\arraybackslash}p{#1}}
\newcommand{\repo}[1]{\path{#1}}
\titleformat{\section}{\normalfont\bfseries\large}{\thesection}{0.5em}{}
\titleformat{\subsection}{\normalfont\itshape\normalsize}{\thesubsection}{0.5em}{}
\captionsetup{font=small,labelfont=bf}
\urlstyle{same}

\begin{document}

\begin{center}
{\small UDC 004.021\par}
\vspace{0.35em}
{\bfseries\itshape Dias Kazikhanov,} Master's student, School of Software Engineering, Astana IT University, Astana, Kazakhstan\par
{\itshape Aigul Meirmanova,} PhD, Postdoctoral Researcher, Assistant Professor, School of Software Engineering, Astana IT University, Astana, Kazakhstan\par
\vspace{0.7em}
{\bfseries\large DEVELOPMENT OF A MULTITASK TRANSFORMER-BASED MODEL AND WEB SYSTEM FOR JOINT SENTIMENT ANALYSIS AND AI-GENERATED REVIEW PROVENANCE CLASSIFICATION: A CONTROLLED MIXED-DOMAIN STUDY\par}
\end{center}

\begin{abstract}
Online reviews combine an evaluative signal with information about how a text was produced. This article presents the completed experimental stage of a dissertation project that develops a multitask Transformer-based model and the ReviewGuard web system for joint review analysis. The model predicts three-way sentiment and a provenance label that is operationalized as human-authored versus AI-generated text in the MAiDE-up resource; this definition is narrower than universal fake-review detection. A corpus of 30,000 records was assembled from four open datasets. TF--IDF with logistic regression, two single-task XLM-RoBERTa models, and one shared-encoder multitask model were compared under one leakage-controlled split and five training seeds. The multitask model reached test Macro-F1 of 0.7709 $\pm$ 0.0087 for sentiment and 0.9526 $\pm$ 0.0094 for provenance, compared with 0.7161 and 0.8998 for the classical baseline. Differences from the single-task Transformer were inconclusive: $-0.0013$ for sentiment and $-0.0046$ for provenance. The result supports a compact shared service, but not a claim of positive transfer or universal authenticity detection. The web layer exposes predictions, confidence, truncation status, and research-context warnings for human-in-the-loop use.
\end{abstract}

\noindent\textbf{Keywords:} sentiment analysis; AI-generated review provenance; provenance classification; multitask learning; XLM-RoBERTa; e-commerce; reproducibility; web system.

\section{Introduction}

Online reviews influence product ranking, merchant reputation, customer
trust, and purchase decisions. An e-commerce analytics service therefore
needs to answer two different questions: what evaluation a reviewer
expresses, and whether the text belongs to the provenance category
represented in the training data. These questions are related because
they use the same review text, but they are not interchangeable. A
negative review is not automatically deceptive, and a classifier score
is not independent evidence of misconduct.

The dissertation topic is to develop a multitask Transformer-based model
and web system for joint sentiment analysis and review authenticity
detection in e-commerce. In the executed study, authenticity is defined
narrowly and reproducibly through the binary MAiDE-up label:
human-authored hotel review versus AI-generated hotel review. The paper
therefore uses provenance classification when describing the empirical
target and reserves broader authenticity language for the application
goal.

Multi-task learning provides a direct way to test whether one encoder
can support both tasks through shared representations. XLM-RoBERTa is
appropriate for the multilingual corpus because it was pretrained across
100 languages \citep{conneau2020xlmr}. However, shared parameters do not guarantee
positive transfer. The comparison with single-task models is essential,
and the training protocol must be documented well enough to separate an
architectural effect from differences in sampling or optimization.

The contribution is fourfold. First, the study reports a completed
30,000-record, five-seed evaluation. Second, it makes label provenance
and the split audit explicit. Third, it separates the scientific
comparison from the supporting web-system implementation. Fourth, it
constrains the interpretation to the tested domain, label definition,
and fixed evaluation protocol.

\section{Related Work}

The multi-task learning formulation treats a shared representation as an
inductive bias that can improve generalization when tasks are related
\citep{caruana1997multitask}. In modern NLP, pretrained Transformer encoders provide the
shared representation, while task-specific heads allow different label
spaces. The benefit is empirical rather than automatic: task imbalance,
label sparsity, and gradient conflict can reduce or reverse transfer.

XLM-RoBERTa extends this setting to multilingual text through
cross-lingual pretraining \citep{conneau2020xlmr}. Its use in this study is motivated by
the corpus composition, not by an assumption that all language slices
are equally represented. The Russian portion dominates the merged
corpus, while the provenance labels come from a multilingual hospitality
dataset. This makes source- and language-aware evaluation important.

Review authenticity has several non-equivalent meanings. Classical
deceptive-opinion studies distinguish truthful and deceptive writing,
whereas newer datasets focus on AI-generated reviews. MAiDE-up is
valuable for the latter setting, but its human-versus-AI labels should
not be generalized to all spam, incentive abuse, account takeover, or
coordinated manipulation \citep{ignat2025maide}. The present study adopts that narrower
semantics.

Finally, a fixed test set and repeated seeds require careful
interpretation. Statistical testing guidance for NLP recommends matching
the test to the experimental unit and reporting uncertainty rather than
treating one p-value as a complete argument \citep{dror2018hitchhiker}. This article
therefore reports seed variation, paired confidence intervals, and
approximate randomization comparisons.

\section{Research Questions and Scope}

The study investigates three questions:

\begin{quote}
• RQ1. Does multitask XLM-RoBERTa improve Macro-F1 over TF-IDF with
logistic regression on sentiment or provenance under the fixed test
protocol?

• RQ2. Does multitask XLM-RoBERTa differ from single-task XLM-RoBERTa
trained for the same task?

• RQ3. Which data, label, and training properties limit transfer beyond
this protocol and the ReviewGuard demonstration system?
\end{quote}

The principal hypotheses are deliberately modest. H1 predicts that
multitask training will not materially underperform single-task training
on sentiment. H2 predicts that both Transformer variants will exceed the
classical baseline on provenance. H3 treats the web interface as an
engineering deliverable, not as evidence of model interpretability or a
substitute for external validation.

\section{Data and Label Provenance}

The processed input contains 30,000 normalized records and is identified
by SHA-256 digest
a296cf7bb8cb0dbb3eebd8a0cc9c478422e9f5959836bc4d20915f9f81c9c5fc. Four
open sources were merged with explicit source identifiers. The source
roles and label semantics are shown in Table 1 \citep{ignat2025maide,rureviews,perekrestok,wildberries}.

\emph{Table 1 -- Corpus composition and label provenance.}

{\def\LTcaptype{none} % do not increment counter
\begin{longtable}[]{@{}
  >{\raggedright\arraybackslash}p{(\linewidth - 6\tabcolsep) * \real{0.2500}}
  >{\centering\arraybackslash}p{(\linewidth - 6\tabcolsep) * \real{0.2500}}
  >{\raggedright\arraybackslash}p{(\linewidth - 6\tabcolsep) * \real{0.2500}}
  >{\raggedright\arraybackslash}p{(\linewidth - 6\tabcolsep) * \real{0.2500}}@{}}
\toprule\noalign{}
\begin{minipage}[b]{\linewidth}\centering
\textbf{Source}
\end{minipage} & \begin{minipage}[b]{\linewidth}\centering
\textbf{n}
\end{minipage} & \begin{minipage}[b]{\linewidth}\centering
\textbf{Sentiment}
\end{minipage} & \begin{minipage}[b]{\linewidth}\centering
\textbf{Provenance}
\end{minipage} \\
\midrule\noalign{}
\endhead
\bottomrule\noalign{}
\endlastfoot
MAiDE-up & 7,000 & Source label; negative/positive & 3,500
human-authored and 3,500 AI-generated \\
Perekrestok & 10,000 & Rating mapping: 1-2 negative, 3 neutral, 4-5
positive & No label \\
RuReviews & 3,000 & Native three-class dataset label & No label \\
Wildberries & 10,000 & Rating mapping: 1-2 negative, 3 neutral, 4-5
positive & No label \\
\end{longtable}
}

All records have a sentiment label. There are 23,734 Russian records and
6,266 records in Turkish, French, Korean, English, Chinese, Spanish,
Italian, German, or Romanian. The domain mixture is 23,000 e-commerce
records and 7,000 hospitality records. Provenance labels are available
only for the 7,000 MAiDE-up records, so e-commerce rows contribute to
sentiment loss but not to provenance loss.

The rating-derived labels are a practical volume compromise, not
independent manual annotation. The sentiment task therefore measures
agreement with the project' s three-class rating mapping.
It should not be interpreted as a universal linguistic definition of
polarity. Similarly, the provenance task measures the MAiDE-up
human-versus-AI distinction and should not be reported as a detector for
every kind of fake review.

\section{Methodology}

\subsection{Leakage-Controlled Split}

Using split seed 42, the corpus was divided into 23,975 training, 3,013
validation, and 3,012 test records. Texts were case-folded and
whitespace-normalized before splitting, then grouped by normalized text
and stratified by source and available-label pair. The audit found 521
exact and 564 normalized duplicate rows in the merged corpus, but zero
cross-partition overlap by record identifier, exact text, or normalized
text for every train, validation, and test pair. The test sentiment
distribution is 710 negative, 204 neutral, and 2,098 positive; the
provenance test is balanced at 350 authentic and 350 AI-generated
examples.

\subsection{Model Architecture}

The classical reference uses TF-IDF unigrams and bigrams with up to
50,000 features and a class-balanced logistic-regression classifier for
each task. The Transformer systems use FacebookAI/xlm-roberta-base. The
shared encoder produces a contextual representation from the first
token. One MLP head predicts three sentiment classes and a second
predicts two provenance classes; each head uses a GELU hidden layer and
dropout of 0.1. For a record with both labels, the loss is the weighted
sum of the two cross-entropies. When one label is absent, only the
available task contributes to the loss. Single-task models use the same
encoder and head design but optimize one label only.

\subsection{Training Protocol}

The local Apple Silicon configuration uses maximum sequence length 192,
batch size 2, four-step gradient accumulation, gradient checkpointing,
AdamW with learning rate 2e-5, weight decay 0.01, 10 percent linear
warmup, a five-epoch maximum, and early stopping with patience 2.
Multi-task training uses a source-balanced sampler and class-weighted
cross-entropy; the single-task loop does not use the source-balanced
sampler. Consequently, the multitask-versus-single-task comparison
evaluates two configured training systems rather than one isolated
architectural switch.

Training was run on a Mac mini M2 Pro with 16 GB unified memory through
the MPS backend. The recorded environment was Python 3.12.12 on macOS
26.2 arm64. Transformer systems used seeds 11, 21, 42, 84, and 126.
Checkpoints were selected using validation Macro-F1 and the test set was
reserved for final evaluation. Macro-F1 is the primary metric because
class balance differs across tasks; accuracy is reported as a secondary
metric.

\subsection{ReviewGuard Web System}

The model is embedded in ReviewGuard, a reproducibility-oriented service
with five stages: data normalization, training and checkpoint export,
checkpoint loading, FastAPI inference, and a browser interface. The
/health endpoint reports service and model readiness; /research-context
returns the training scope and dataset warnings; and POST /analyze
accepts a review and returns sentiment and provenance labels with
confidence scores.

The response also contains a lightweight transparency object: ranked
task probabilities, token count, configured maximum length, a truncation
flag, prediction margins, risk flags, and research-context warnings.
This is an operational aid for a human reviewer, not a causal
explanation of the model and not evidence that the predicted provenance
label is factually true.

\section{Experimental Design}

Each Transformer variant was trained for five random seeds on the same
fixed split. The classical baseline is deterministic under the fixed
configuration. For Transformer systems, the article reports mean and
standard deviation over seeds and a 95 percent t-interval for the mean.
Paired comparisons use the same test examples: 2,000 bootstrap samples
produce an interval for the seed-averaged difference, and an approximate
randomization test uses 2,000 permutations. These intervals describe
uncertainty within the chosen split; they do not quantify uncertainty
from new sources, domains, or languages. We also computed worst-slice
Macro-F1 by language, source, and domain with a minimum slice support of
25.

The comparison is intentionally bounded. It tests a shared encoder
against two single-task encoders and a shallow text baseline under one
data construction. It does not claim that the chosen source quotas
represent the population of e-commerce reviews, that MAiDE-up labels
cover human deception, or that a production threshold has been
calibrated.

\section{Results}

\subsection{Main Test Metrics}

For class $c$, precision is $P_c=TP_c/(TP_c+FP_c)$ and recall is $R_c=TP_c/(TP_c+FN_c)$. The reported Macro-P, Macro-R, and Macro-F1 are unweighted means over classes; Weighted-F1 weights class F1 by test-set support. All metrics are computed on the fixed test partition.

\emph{Table 2 -- Test metrics. Transformer values are mean +/- standard
deviation over five seeds.}

{\def\LTcaptype{none} % do not increment counter
\begin{longtable}[]{@{}
  >{\raggedright\arraybackslash}p{(\linewidth - 6\tabcolsep) * \real{0.2500}}
  >{\centering\arraybackslash}p{(\linewidth - 6\tabcolsep) * \real{0.2500}}
  >{\raggedright\arraybackslash}p{(\linewidth - 6\tabcolsep) * \real{0.2500}}
  >{\raggedright\arraybackslash}p{(\linewidth - 6\tabcolsep) * \real{0.2500}}@{}}
\toprule\noalign{}
\begin{minipage}[b]{\linewidth}\centering
\textbf{Model family}
\end{minipage} & \begin{minipage}[b]{\linewidth}\centering
\textbf{Task}
\end{minipage} & \begin{minipage}[b]{\linewidth}\centering
\textbf{Macro-F1}
\end{minipage} & \begin{minipage}[b]{\linewidth}\centering
\textbf{Accuracy}
\end{minipage} \\
\midrule\noalign{}
\endhead
\bottomrule\noalign{}
\endlastfoot
TF-IDF + logistic regression & Sentiment & 0.7161 & 0.8320 \\
Single-task XLM-R & Sentiment & 0.7721 +/- 0.0087 & 0.8958 +/- 0.0038 \\
Multitask XLM-R & Sentiment & 0.7709 +/- 0.0087 & 0.8950 +/- 0.0047 \\
TF-IDF + logistic regression & Provenance & 0.8998 & 0.9000 \\
Single-task XLM-R & Provenance & 0.9571 +/- 0.0088 & 0.9571 +/-
0.0087 \\
Multitask XLM-R & Provenance & 0.9526 +/- 0.0094 & 0.9526 +/- 0.0094 \\
Majority baseline & Sentiment & 0.2737 & 0.6965 \\
Majority baseline & Provenance & 0.3333 & 0.5000 \\
\end{longtable}
}

Multitask XLM-R exceeds the TF-IDF baseline by 0.0548 Macro-F1 for
sentiment and 0.0528 for provenance. The single-task Transformer has
slightly higher means on both tasks, but the gaps are small relative to
seed variation. The main practical result is therefore a compact shared
model with comparable average quality, not a demonstrated accuracy
advantage over separate models.

\subsection{Paired Comparisons}

\emph{Table 3 -- Paired Macro-F1 comparisons for the multitask model.
Delta equals multitask minus comparator.}

{\def\LTcaptype{none} % do not increment counter
\begin{longtable}[]{@{}
  >{\raggedright\arraybackslash}p{(\linewidth - 8\tabcolsep) * \real{0.2000}}
  >{\centering\arraybackslash}p{(\linewidth - 8\tabcolsep) * \real{0.2000}}
  >{\centering\arraybackslash}p{(\linewidth - 8\tabcolsep) * \real{0.2000}}
  >{\raggedright\arraybackslash}p{(\linewidth - 8\tabcolsep) * \real{0.2000}}
  >{\centering\arraybackslash}p{(\linewidth - 8\tabcolsep) * \real{0.2000}}@{}}
\toprule\noalign{}
\begin{minipage}[b]{\linewidth}\centering
\textbf{Task}
\end{minipage} & \begin{minipage}[b]{\linewidth}\centering
\textbf{Comparator}
\end{minipage} & \begin{minipage}[b]{\linewidth}\centering
\textbf{Delta}
\end{minipage} & \begin{minipage}[b]{\linewidth}\centering
\textbf{95\% CI}
\end{minipage} & \begin{minipage}[b]{\linewidth}\centering
\textbf{p (approx. randomization)}
\end{minipage} \\
\midrule\noalign{}
\endhead
\bottomrule\noalign{}
\endlastfoot
Sentiment & Single-task XLM-R & -0.0013 & {[}-0.0105, 0.0081{]} &
0.8016 \\
Sentiment & TF-IDF + logistic regression & +0.0548 & {[}0.0337,
0.0732{]} & 0.0005 \\
Provenance & Single-task XLM-R & -0.0046 & {[}-0.0109, 0.0020{]} &
0.1829 \\
Provenance & TF-IDF + logistic regression & +0.0528 & {[}0.0305,
0.0749{]} & 0.0005 \\
\end{longtable}
}

\emph{Table 4 -- Additional macro-averaged test metrics. Values are mean
+/- standard deviation over five seeds.}

{\def\LTcaptype{none} % do not increment counter
\begin{longtable}[]{@{}
  >{\raggedright\arraybackslash}p{(\linewidth - 8\tabcolsep) * \real{0.2000}}
  >{\raggedright\arraybackslash}p{(\linewidth - 8\tabcolsep) * \real{0.2000}}
  >{\centering\arraybackslash}p{(\linewidth - 8\tabcolsep) * \real{0.2000}}
  >{\centering\arraybackslash}p{(\linewidth - 8\tabcolsep) * \real{0.2000}}
  >{\centering\arraybackslash}p{(\linewidth - 8\tabcolsep) * \real{0.2000}}@{}}
\toprule\noalign{}
\begin{minipage}[b]{\linewidth}\centering
\textbf{Model}
\end{minipage} & \begin{minipage}[b]{\linewidth}\centering
\textbf{Task}
\end{minipage} & \begin{minipage}[b]{\linewidth}\centering
\textbf{Macro-P}
\end{minipage} & \begin{minipage}[b]{\linewidth}\centering
\textbf{Macro-R}
\end{minipage} & \begin{minipage}[b]{\linewidth}\centering
\textbf{Weighted F1}
\end{minipage} \\
\midrule\noalign{}
\endhead
\bottomrule\noalign{}
\endlastfoot
TF-IDF + logistic regression & Sentiment & 0.6919 & 0.7688 & 0.8446 \\
Single-task XLM-R & Sentiment & 0.7740 +/- 0.0028 & 0.7737 +/- 0.0173 &
0.8964 +/- 0.0035 \\
Multitask XLM-R & Sentiment & 0.7687 +/- 0.0104 & 0.7762 +/- 0.0180 &
0.8962 +/- 0.0043 \\
TF-IDF + logistic regression & Provenance & 0.9034 & 0.9000 & 0.8998 \\
Single-task XLM-R & Provenance & 0.9577 +/- 0.0083 & 0.9571 +/- 0.0087 &
0.9571 +/- 0.0088 \\
Multitask XLM-R & Provenance & 0.9531 +/- 0.0090 & 0.9526 +/- 0.0094 &
0.9526 +/- 0.0094 \\
\end{longtable}
}

\emph{Table 5 -- Per-class test metrics for the multitask model. Values
are mean +/- standard deviation over five seeds.}

{\def\LTcaptype{none} % do not increment counter
\begin{longtable}[]{@{}
  >{\raggedright\arraybackslash}p{(\linewidth - 8\tabcolsep) * \real{0.2000}}
  >{\raggedright\arraybackslash}p{(\linewidth - 8\tabcolsep) * \real{0.2000}}
  >{\centering\arraybackslash}p{(\linewidth - 8\tabcolsep) * \real{0.2000}}
  >{\centering\arraybackslash}p{(\linewidth - 8\tabcolsep) * \real{0.2000}}
  >{\centering\arraybackslash}p{(\linewidth - 8\tabcolsep) * \real{0.2000}}@{}}
\toprule\noalign{}
\begin{minipage}[b]{\linewidth}\centering
\textbf{Task}
\end{minipage} & \begin{minipage}[b]{\linewidth}\centering
\textbf{Class}
\end{minipage} & \begin{minipage}[b]{\linewidth}\centering
\textbf{Precision}
\end{minipage} & \begin{minipage}[b]{\linewidth}\centering
\textbf{Recall}
\end{minipage} & \begin{minipage}[b]{\linewidth}\centering
\textbf{F1}
\end{minipage} \\
\midrule\noalign{}
\endhead
\bottomrule\noalign{}
\endlastfoot
Sentiment & Negative & 0.8700 +/- 0.0156 & 0.8566 +/- 0.0192 & 0.8630
+/- 0.0083 \\
Sentiment & Neutral & 0.4887 +/- 0.0415 & 0.5284 +/- 0.0663 & 0.5041 +/-
0.0205 \\
Sentiment & Positive & 0.9475 +/- 0.0063 & 0.9437 +/- 0.0047 & 0.9456
+/- 0.0023 \\
Provenance & Authentic & 0.9654 +/- 0.0073 & 0.9389 +/- 0.0208 & 0.9518
+/- 0.0100 \\
Provenance & AI-generated & 0.9408 +/- 0.0190 & 0.9663 +/- 0.0077 &
0.9533 +/- 0.0088 \\
\end{longtable}
}

The additional metrics expose an important asymmetry that Macro-F1 alone
hides: the neutral sentiment class is substantially weaker (F1 0.5041
+/- 0.0205) than negative and positive sentiment, while the two
provenance classes remain balanced. The 95 percent t-based confidence
intervals across seeds are {[}0.7600, 0.7817{]} for multitask sentiment
Macro-F1 and {[}0.9409, 0.9642{]} for multitask provenance Macro-F1.
Paired Cohen' s dz values are -0.1383 and -0.2724 for
multitask versus single-task sentiment and provenance, respectively, and
6.2744 and 5.6106 versus the classical baseline; because only five seeds
are available, these effect sizes are descriptive. AUROC, PR-AUC,
calibration error, and Brier score are not reported because the release
artifacts contain aggregate confusion-matrix metrics but no per-example
probability outputs; these metrics should be added in the next
evaluation run.

The majority baselines in Table 2 are 0.2737 Macro-F1 and 0.6965
accuracy for sentiment, and 0.3333 Macro-F1 and 0.5000 accuracy for
provenance. For the four Macro-F1 comparisons in Table 3, Holm-adjusted
p-values are 0.0020 for both multitask-versus-baseline comparisons,
0.3658 for provenance versus single-task, and 0.8016 for sentiment
versus single-task. These inferences remain conditional on the fixed
test split and are exploratory because only five training seeds are
available.

A robustness audit across language, source, and domain slices also
limits the generalization claim. Across the five seeds, the
worst-language Macro-F1 averaged 0.5688 +/- 0.0220 for sentiment and
0.8905 +/- 0.0267 for provenance; the worst-source and worst-domain
sentiment Macro-F1 values averaged 0.6210 +/- 0.0039 and 0.6223 +/-
0.0030. Several slices contain no neutral sentiment examples, so these
slice scores are diagnostic and should not be interpreted as fully
comparable three-class estimates.

The intervals for the multitask-versus-single-task comparisons cross
zero, so the experiment does not establish positive transfer. Against
the classical baseline, the intervals are positive and the randomization
p-value is 0.0005 for both tasks, the smallest value resolvable with
2,000 permutations. The answer to RQ1 is positive within the fixed
protocol; the answer to RQ2 is that no reliable difference is
established.

\section{Discussion}

The result pattern is consistent across tasks: Transformer
representations are clearly stronger than the selected TF-IDF reference,
while the distinction between shared and single-task Transformer
training is small. This may indicate that the two targets use partially
shared lexical and semantic information without providing enough
independent provenance supervision to produce measurable positive
transfer. It may also reflect the non-isolated sampler difference. The
current evidence cannot distinguish these explanations.

For the dissertation system, the compact shared encoder remains useful.
It provides one exportable checkpoint, one API contract, and one web
interface for two predictions. That engineering advantage should not be
converted into a scientific claim that the shared model is more
accurate. A deployment should expose confidence and truncation
information, require human review for high-impact decisions, and retain
the source and label warnings returned by the research context.

The web system also clarifies the boundary between a model and a
moderation policy. A provenance score is a research signal tied to
MAiDE-up' s human-versus-AI definition. It is not a
decision that a reviewer is fraudulent, and it should not be used for
automatic penalties without independent validation, calibration, and an
appeal process.

\section{Limitations and Next Steps}

\begin{quote}
• Provenance validity: all provenance labels come from MAiDE-up and the
hotel domain. They contrast human-authored and AI-generated reviews and
do not test human-written deception, incentive abuse, or coordinated
spam.

• Heterogeneous sentiment supervision: Perekrestok and Wildberries
labels are derived from ratings, so source and label-generation effects
may be confounded with language and domain effects.

• One fixed split: five training seeds measure initialization
sensitivity but do not replace source-held-out, language-held-out, or
external evaluation. Before redistribution, source-specific provenance
subtype metadata must also be audited against authenticity\_label; the
reported metrics use the authenticity\_label field directly.

• Duplicate structure and quota sampling: the audit prevents observable
cross-partition text leakage, but normalized duplicates remain inside
the merged corpus and fixed source quotas are not a random sample of all
reviews.

• Non-isolated ablation: the multitask system uses source-balanced
sampling while the single-task system does not. A stricter architectural
ablation must equalize sampling and optimization schedules.

• Deployment validation: probabilities were not calibrated, a moderation
threshold was not selected, and the web system is decision support
rather than an autonomous enforcement service.
\end{quote}

The next dissertation stage should collect or obtain independent
provenance labels in an e-commerce domain, repeat the comparison with
identical samplers, report source- and language-held-out results,
calibrate probabilities, and evaluate the complete API workflow with
representative user scenarios. These steps would test generalization
rather than only extending the training corpus.

\section{Conclusion}

This article presents a reproducible multitask XLM-RoBERTa model and
ReviewGuard web system for joint sentiment analysis and review
provenance classification. On a leakage-controlled, partially labelled
corpus of 30,000 records, multitask training increased Macro-F1 over
TF-IDF with logistic regression from 0.7161 to 0.7709 for sentiment and
from 0.8998 to 0.9526 for provenance. Single-task Transformer systems
had slightly higher means, and paired intervals did not confirm a
difference.

The defensible conclusion is therefore conditional. A shared encoder can
support both predictions in one service without a clear average-quality
loss under this protocol, but the experiment does not prove positive
transfer and does not establish universal fake-review detection. The
strongest next test is an external, source-held-out and domain-held-out
evaluation with independently defined provenance labels.

\section*{Data Availability and Reproducibility}

The project repository contains the data-normalization code, training
CLI, configuration snapshot, exported model format, API service, and
aggregated reports. The principal reproducibility artifacts include
data/processed/joint\_reviews.article30k.jsonl,
configs/model.article30k.mac\_m2\_16gb.yaml, the multiseed report and
statistics directories, and the split/leakage audit. The complete-run
manifest records commit 64cf0a580d4a28434454aba763bec51ff052c689, Python
3.12.12, macOS 26.2 arm64, the MPS backend, split seed 42, and training
seeds 11, 21, 42, 84, and 126. The corpus digest is given in Section 4.
Redistribution of source records follows the original dataset terms;
exact reproduction should use the cited dataset cards and recorded
configuration rather than silently substituting newer versions. The
repository was dirty during the run, so the modified paths must be
preserved with any supplementary artifact.

\section*{Ethics Statement}

No participants were recruited and no new personal data were collected.
The study uses textual records from open research sources. The release
package does not reproduce raw review text or direct identifiers; the
locally processed corpus may retain source-specific identifiers for
audit purposes and must not be redistributed without checking the
original licenses and applying an explicit privacy/PII review.
Provenance predictions are research signals, not standalone evidence of
author misconduct.

\section*{CRediT Author Contributions}

Dias Kazikhanov: conceptualization, methodology, software, data
curation, formal analysis, validation, visualization, writing - original
draft, and writing - review and editing. Supervisor: A. Meirmanova.

\section*{Funding and Conflict of Interest}

No external funding or conflict-of-interest declaration was available in
the project materials at the time of writing; the authors should confirm
the official statement required by the target venue before submission.

\section*{AI Use Disclosure}

Generative AI was used as an auxiliary tool for manuscript structure and
language drafting. Numerical results, experimental configuration, data
claims, and references were checked against repository artifacts and the
cited primary sources. The author remains responsible for the final
manuscript, interpretation, and compliance with the submission
venue' s policies.

\bibliographystyle{unsrtnat}
\bibliography{article_improved_en_meirmanova_reviewed}

\section*{Information about Authors}

\emph{Dias Kazikhanov -- Master's student, School of Software
Engineering, Astana IT University, Astana, Kazakhstan.}

\emph{Aigul Meirmanova -- PhD, Postdoctoral Researcher, Assistant
Professor, School of Software Engineering, Astana IT University, Astana,
Kazakhstan.}

\end{document}
